Viterbo's conjecture states that, among convex bodies of fixed volume, the
ball has the largest symplectic capacity. Haim--Kislev and Ostrover
disproved it with a Lagrangian product of two regular pentagons, the
Haim--Kislev--Ostrover (HKO) body. We prove that this body locally maximizes
the systolic ratio among polytopes with exactly ten facets. Among bodies
sufficiently near HKO, equality in this bound holds precisely for its
translations, positive dilations and linear symplectic images. We exhibit a
second ten-facet local maximum: a rational hexagon--quadrilateral product of
ratio $1$, with a two-parameter local equality family modulo the same
symmetries; it is not a new counterexample.

For products of two independently linearly deformed regular pentagons we
prove an explicit capacity formula and the sharp bound $(3+\sqrt5)/5$, so no
deformation within this family improves the HKO ratio. Every symplectic
product of two full-dimensional planar convex bodies admits a
capacity-increasing rigid rotation, while the Lagrangian endpoint can have
smaller capacity. The proofs rest on finite capacity formulas and their
first-order variation. The computations use an implementation of
Haim--Kislev's quadratic program and, as an independent check, a flow-graph
algorithm built on the local passage geometry of Chaidez and Hutchings,
which is correct under explicit regularity assumptions.

In sampled populations of four-dimensional convex polytopes, the total
symplectic area of the two-dimensional faces, normalized by volume, has a
strong inverse association with the computed systolic ratio. The association
weakens among low-area or high-ratio subsets, and exact laws for rotated
regular pentagon and equilateral-triangle products order area and ratio in
opposite ways. None of the exploratory searches produced a numerical ratio
above one.
